% Worked Examples: Concept 2.6.4 - Bode's Relations and Minimum Phase Systems
\documentclass[11pt,a4paper]{article}
\usepackage{worked-examples-template}

\title{\textbf{Worked Examples}\\[0.5em]
\large Concept 2.6.4: Bode's Relations and Minimum Phase Systems\\[0.3em]
\normalsize \coursesubtitle{} -- \coursetitle}
\author{Assoc.\ Prof.\ William Robertson \and Dr Sean McGowan}
\date{\today}

\begin{document}

\maketitle
\tableofcontents
\newpage

\section{Concept 2.6.4: Minimum Phase Systems}

\subsection{Example 1: Identifying Minimum Phase Systems}

\paragraph{Problem:} Determine which of the following transfer functions are minimum phase systems:

(a) $G_1(s) = \frac{s+2}{(s+1)(s+3)}$

(b) $G_2(s) = \frac{s-2}{(s+1)(s+3)}$

(c) $G_3(s) = \frac{2e^{-0.5s}}{s+1}$

For each system, explain your reasoning.

\paragraph{Solution:}

A system is \textbf{minimum phase} if all its poles and zeros are in the left half of the complex plane (i.e., have negative real parts), and it has no time delays.

\textbf{Part (a): $G_1(s) = \frac{s+2}{(s+1)(s+3)}$}

Zeros: $s = -2$ (left half-plane) ✓

Poles: $s = -1, -3$ (both left half-plane) ✓

No time delay ✓

\textbf{Result:} $G_1(s)$ is a \textbf{minimum phase system}.

\textbf{Part (b): $G_2(s) = \frac{s-2}{(s+1)(s+3)}$}

Zeros: $s = +2$ (right half-plane) ✗

Poles: $s = -1, -3$ (both left half-plane) ✓

No time delay ✓

\textbf{Result:} $G_2(s)$ is \textbf{NOT minimum phase} because it has a right half-plane zero (at $s = +2$). This is called a \textit{non-minimum phase zero}.

\textbf{Part (c): $G_3(s) = \frac{2e^{-0.5s}}{s+1}$}

Zeros: None (the time delay $e^{-0.5s}$ is not a zero in the traditional sense)

Poles: $s = -1$ (left half-plane) ✓

Has time delay: $e^{-0.5s}$ ✗

\textbf{Result:} $G_3(s)$ is \textbf{NOT minimum phase} because it contains a time delay. Time delays introduce infinite phase lag without affecting magnitude, making the system non-minimum phase.

\textbf{Summary Table:}

\begin{center}
\begin{tabular}{|c|c|c|c|}
\hline
System & LHP Zeros & LHP Poles & No Delay & Minimum Phase? \\
\hline
$G_1(s)$ & ✓ & ✓ & ✓ & Yes \\
$G_2(s)$ & ✗ (RHP zero) & ✓ & ✓ & No \\
$G_3(s)$ & ✓ & ✓ & ✗ (has delay) & No \\
\hline
\end{tabular}
\end{center}

\textbf{Key Properties of Minimum Phase Systems:}
\begin{itemize}
\item Unique relationship between magnitude and phase response
\item Phase can be reconstructed from magnitude (and vice versa)
\item Minimum possible phase lag for a given magnitude response
\item Generally easier to control than non-minimum phase systems
\end{itemize}

\textbf{Practical Implications:}
\begin{itemize}
\item Non-minimum phase zeros (RHP zeros) impose fundamental limitations on achievable bandwidth and performance
\item Time delays make control more challenging and limit achievable phase margin
\item Example: inverted pendulum has non-minimum phase dynamics
\end{itemize}
\end{document}
